\chapter{The nonlinear renormalization step}
\label{ch09:renormalization}

A block transformation acts on the density, not on its chosen list
of local terms. Our task is to express its exact output in the
coordinates of Chapter~\ref{ch04:blocking}, while improving the
bounds enough to repeat the operation. The organization by local
interactions and large-field regions follows the constructive
renormalization method of Ba{\l}aban
\cite{Balaban1988Expansions}. We use the cluster theorem of
Chapter~\ref{ch08:chapter}; the work here is the identification,
normalization and estimation of its concrete activities.

Fix one adjacent pair of scales, with fine scale $m+1$ and coarse
scale $m$. Write $L=L_m$, $B=B_m$, $t=t_m$, and
$\mathfrak b=\widehat\beta_{m+1}$. The actual incoming inverse coupling is
$\beta'\in[\mathfrak b/2,2\mathfrak b]$.
All event indicators use the fixed reference thresholds of
\eqref{ch04:thresholds}. They do not move when $\beta'$ varies.
Write the prescribed coarse path action as
\[
 P^{\rm path}_{N,m}(\beta,U)=\beta\mathcal K(U)+\mathcal L(U),
\]
where $\mathcal K$ includes its cubic and quartic terms and
$\mathcal L$ is the quadratic loop correction.
This identity defines $\mathcal K,\mathcal L$ from
\eqref{ch04:path-action}. Neither is the unknown effective action.

For this induction the derivative orders and spatial exponents
are stronger than the minimum allowed in the density definition:
take $q_{m+1}\ge q_m\ge6$ and
$\varpi_{m+1}\ge\varpi_m\ge2$, with the prescribed derivative
reserves for the local integral estimates. In particular the
normalized affine directions used below belong to the activity
seminorm. At the finest scale choose $0<c_{\rm dir}\le1$.

The estimates below distinguish proved algebraic implications
from the analytic inputs of Chapters~\ref{ch06:large-fields}
and~\ref{ch07:matching}. In particular a product estimate under
one Gaussian law, a local low-order matching estimate, and a
one-loop drift estimate are different inputs.

\section{The exact local integral and its first expansion}
\label{ch09:integral}

Let
$\mathscr Z_{\rm blk}=\mathcal B_m^{\beta'}\operatorname{dens}(\rho')$,
with the exact kernel \eqref{ch04:block-operator}.
In each sector one first performs the local coordinate change
on its active links and leaves the other links in probability
Haar measure. Chapter~\ref{ch06:large-fields} gives the exact
Gaussian identity, including the difference between the chosen
finite-range covariance and the inverse precision.

Here is the scalar bookkeeping in a form useful for checking
reassembly. Suppose there are $N_{\rm free}$ free tree-gauge
bonds, of which $n_{\rm act}$ use Lie algebra coordinates, and
let $D$ be the dimension after dummy Gaussian coordinates have
been added. Set $g=(\beta')^{-1/2}$, $q=\dim\gLie$, and
$c_{\rm H}$ equal to the density of probability Haar measure
in the chosen exponential coordinates at zero.
With $a=\mu+gz$, the measure multiplier before the Haar
logarithm is
\[
 (2\pi)^{-(D-n_{\rm act}q)/2}g^{n_{\rm act}q}c_{\rm H}^{n_{\rm act}}
 =(2\pi)^{-D/2}\mathfrak m^{N_{\rm free}}
                  e^{-(N_{\rm free}-n_{\rm act})\log\mathfrak m},
 \quad
 \mathfrak m=c_{\rm H}(2\pi g^2)^{q/2}.
\]
The dummy Gaussian integral equals one after the inactive Haar
variables have been frozen; Fubini then permits their integration.
The flat reference determinant uses the same forest as the
sector. It is extracted once, not separately for every eventual
connected component. The local precision correction, Haar
logarithm, missing-coordinate factor and stationary residual
remain among the local exponential factors.
Keeping them is what recovers the original Haar integral.

Expand exponentials as $e^V=1+(e^V-1)$ and retain all sector
allocation indicators. A complete dependence region contains
every Haar coordinate, tag variable and Gaussian covariance
halo read by its factor. If two such supports are disjoint,
the Haar and tag coordinates are independent; the Gaussian
cross covariance is zero because the chosen covariance has
finite range. Their integrals therefore factor.
Connected components of intersecting dependence regions define
the first polymer activities.

This factorization concerns dependence regions. It does not
assert that large-field activities with disjoint exclusion
supports $Y_i$ have disjoint full dependence regions.
The large-field indicators in their eligibility criteria can read outside $Y_i$.
Those conditions remain in every exact-cover expression.

The required simultaneous estimate for a connected product is
\begin{equation}
 \int G_M\,d\nu\le
 P_j(\mathfrak b)(1+s_M)^{d_j}
       e^{-r_j\beta't_{m+1}^2|A_M|}
       \prod_{a\in M}u_a .
 \label{ch09:joint-input}
\end{equation}
Here $\nu$ is the one sector Gaussian/Haar law,
$G_M$ majorizes the product and its prescribed ordered
derivatives, $s_M$ counts the full dependence region in
radius-$B$ blocks, and $A_M$ is its selected large-field condition set.
The function $P_j$ has subpower growth:
$\log P_j(\mathfrak b)/\log\mathfrak b\to0$ at fixed stage and
ratio. Distinct cutoff-failure labels each contribute their
own factor to the product. The bounded-jet estimate
\eqref{ch06:actual-jet-integral}, applied to the primitive factors
constructed in Chapter~\ref{ch06:large-fields}, proves this joint
bound on their common law. The alternative quadratic-tilt criterion
in Proposition~\ref{ch06:common-majorant} is not required here.
Separate bounds on $\int|F_a|\,d\nu$ cannot
replace this estimate.

The elementary power counting is as follows. The entries include
the prescribed finite number of derivatives.
\[
\begin{array}{l|l}
\text{factor after deterministic subtraction}&\text{majorant}\\ \hline
\text{stationary cubic remainder}&P_j\beta'g^3\\
\text{centered determinant variation}&P_j(t_{m+1}+g)\\
\text{Haar remainder}&P_j(t_{m+1}+g)^2\\
\text{inherited quadratic-loop variation}
   &P_j\Lambda_{m+1}^{\rm loop}(t_{m+1}+g)g\\
\text{covariance and stationary residual}
   &P_j\mathfrak b^{c_j}e^{-c'_jB}\\
\text{incoming regular activity}&P_j\epsilon_{m+1}^{\rm reg}\\
\text{cutoff failure at a distinct label}&P_j\mathfrak b^{-1/2}
\end{array}
\]
For example $\beta'g^3=g$, and the stationary mean
is $O(gM_{m+1}^{\rm cut})$, so a cubic Taylor term is bounded
by $g$ times a polynomial in $M_{m+1}^{\rm cut}$ and Gaussian
coordinates. The latter is estimated inside the common integral.
The other rows follow from the determinant inverse identity,
the second-order vanishing of the Haar correction, and
the window errors of Chapter~\ref{ch06:large-fields}.
The loop coefficient bound must grow subpolynomially at fixed ratio.

After increasing the inverse coupling threshold these rows are bounded
by $\mathfrak b^{-7/16}$, since a subpower factor can be bounded
by $\mathfrak b^{1/16}$. Read-count polynomials are absorbed by
additional exponential support decay. The first component expansion
and then Theorem~\ref{ch08:kp} give an absolutely convergent
logarithm of the ordinary gas without large-field sites. Components attached to
a prescribed large-field site set remain in the exact-cover sum.
The deletion identity \eqref{ch08:deletion-identity} accounts
for ordinary polymers excluded by a selected large-field region.
Repeated polymer labels in the logarithm are counted at every
occurrence, with the factor $1/n!$ of \eqref{ch08:log}.

All rearrangements can first be made for finite selections.
Under \eqref{ch09:joint-input} and the strict cluster criterion,
the absolute envelopes allow passage to the countable sums.
Thus this procedure gives an identity for $\mathscr Z_{\rm blk}$,
rather than an asymptotic expansion for its logarithm.

\section{Improving the regular remainder}
\label{ch09:improvement}

The first bound $\mathfrak b^{-7/16}$ is too large for the
target $\epsilon_m^{\rm reg}=\widehat\beta_m^{-11/10}$.
There are three different sources of improvement.

\subsection{Terms containing several elementary factors}

Give every elementary ordinary factor a degree one, and give
a first-expansion polymer the sum of its factors' degrees.
A cluster degree is the sum over its polymer occurrences.
Multiply each elementary majorant by
$h=\mathfrak b^{3/8}$ before performing either expansion.
The resulting local smallness parameter is
$h\mathfrak b^{-7/16}=\mathfrak b^{-1/16}$.
The weighted rooted bound of
Section~\ref{ch08:tail}, applied with these inflated activities,
bounds the whole sum by $P_jh\mathfrak b^{-7/16}$.
Restoring the original weights for terms of total degree at
least three gives
\begin{equation}
 P_jh^{-3}(h\mathfrak b^{-7/16})
       =P_j\mathfrak b^{-19/16}
       \le\mathfrak b^{-23/20}.
 \label{ch09:degree-gain}
\end{equation}
The gap $19/16-23/20=3/80$ absorbs the remaining subpower factor.
This counts elementary factors, not just vertices in the
second cluster expansion; a single polymer may already have
degree three. The same calculation applies after placing a
fixed number of derivatives.

An incoming regular activity accompanied by a further factor
is bounded by the rooted estimate with that incoming activity
distinguished:
$P_j\epsilon_{m+1}^{\rm reg}\mathfrak b^{-1/2}$.
Long path terms have an extra factor $e^{-c_jB}$.
Both fit the right side of \eqref{ch09:degree-gain} after
increasing the threshold.

\subsection{The linear inherited regular term}

The component linear in an incoming regular activity must use
its affine-Hessian cancellation. In local radial gauge, the
lifted potential splits into parts with inverse-ratio sizes
$L^{-1},L^{-2},L^{-3}$: a constant part, its affine variation,
and its remainder. The constant part is removed by a local
gauge transformation. Curvature therefore begins at weight
two. Invariance under constant gauge transformations remove linear curvature terms,
and the incoming affine-Hessian condition removes the
quadratic curvature term of weight four. Every surviving
selection in the ordered Taylor expansion has weight at least
five.

Here is the support estimate needed to sum that local
calculation. Let $d$ be the incoming covering number.
For $d^3<L$ the inspected support admits a nonwrapping
radial gauge and the per-label gain is $L^{-5}$.
For the other supports use the original bound. Both cases
are bounded by
\[
 \delta_L(d)\le L^{-5}(1+d)^{15},\qquad
 \delta_L(d)=
 \begin{cases}L^{-5},&d^3<L,\\1,&d^3\ge L.\end{cases}
\]
The full chain rule adds a fixed polynomial in $d$.
Choose the scale ratio and incoming support exponent so that, after
support compression, the outgoing weight uses at most half of the
incoming exponential decay. The remaining decay absorbs the polynomial
in $d$.
There are $L^4$ fine anchors above each coarse anchor.
Consequently the inherited principal norm is at most
$V_jL^{-1}\epsilon_{m+1}^{\rm reg}$, with $V_j$ fixed before
$L$.

For the comparison with the actual integral use the same forest
and physical local coordinates on both sides. The estimates of
Section~\ref{ch06:linear-comparison} give a covariance difference
$O(P(B)t_f)$ and a mean difference, in Gaussian units, of order
$P(B)g(M_f^{\rm cut})^2+P(B)M_f^{\rm cut}e^{-cB}$,
where $t_f=t_{m+1}$ and $g=(\beta')^{-1/2}$. Each is
$\mathfrak b^{-1/2+o(1)}$. To see where the gain occurs, for
a covariance interpolation $C_s$ with fixed positive lower bound,
the logarithmic density derivative is
\[
 \frac12\left\{z^*C_s^{-1}\dot C_s C_s^{-1}z
                      -\tr(C_s^{-1}\dot C_s)\right\}.
\]
A varying mean adds its linear Gaussian score. Every coefficient
contains a covariance or normalized-mean difference. Its fixed
moments therefore contain the small factor just stated, times a
polynomial in the local dimension.

Hold the physical links in the incoming activity fixed by the
protected coordinate change of Chapter~\ref{ch06:large-fields}.
The nonlinear coordinate correction is then differentiated through
its Jacobian and density, where the same gain is available.
This avoids requiring one extra derivative of the incoming activity
when adjacent levels have equal derivative orders. Telescope all
remaining factor products, marking one changed factor. The
nonquadratic local weights contribute a gain of order $g$ with
subpower losses; a second incoming activity contributes
$O((\epsilon_{m+1}^{\rm reg})^2)$ instead. Common Gaussian moments
and the unused support exponent sum the resulting polynomial costs.
Together with the principal estimate this proves
\[
 \mathcal E_{\sigma+3}(\text{inherited term})
 \le \left(V_jL^{-1}+
       \omega_{j,L}(\widehat\beta_m)\right)
                    \epsilon_{m+1}^{\rm reg},
 \qquad
 \omega_{j,L}(\widehat\beta_m)
       =\widehat\beta_m^{-1/2+o(1)}.
\]
For this comparison the local windows and their dependence halos
must be clean. A window meeting a selected large-field component is
assigned to that component; a failed remote cutoff has the joint
Gaussian-tail bound of Chapter~\ref{ch06:large-fields}. Large
remaining supports gain $e^{-cB}$ from the unused incoming support
weight. Thus no comparison between unrelated event indicators or a seeded and
an unseeded law has been made. The actual interpolation estimates
are essential; the scaling-dimension count alone would not prove
the comparison remainder.

\subsection{Low-order terms and localization}

One- and two-factor terms must be combined before estimating them.
The degree-four classical matching and quadratic determinant matching
must apply to the same integral and stationary branch. Rather than
use a formal cumulant remainder, we retain a finite stationary
region and estimate its complement. The resulting centered
auxiliary bound, with $h=\mathfrak b^{-7/16}$, is
\begin{equation}
 P_j(\mathfrak b)\left\{
 \mathfrak b^{-1}h+\mathfrak b h^5+h^3\right\}
                         +\mathfrak b^{-K},
 \label{ch09:aux-input}
\end{equation}
where $K$ can be any fixed positive number. Its largest power is
$\mathfrak b^{-19/16+o(1)}$, which is smaller than
$\mathfrak b^{-23/20}$. The argument below specifies the actual
sector and stationary estimates needed to obtain this bound.

% !TEX root = ../monograph.tex
\subsection{The stationary integral with its cutoff retained}
\label{ch09:stationary-remainder}

We now give the local analytic implication behind the auxiliary bound.
It is an application of Laplace's method, for which we refer to
\cite[Chapter II]{Wong2001}, with three features that must be checked
here: the integration dimension grows with the auxiliary period, the
Hessian at the critical configuration depends on the coarse field, and the integration
cutoff cannot be discarded before estimating its complement.
All actions and matrices in the following argument belong to the same
\emph{pure} small-field block integral, the block integral of the path action alone: incoming activities are omitted, the
prescribed path action and block kernel are unchanged, and the fine-field
indicator $\mathbf1_{\{\chi(U_p)\le t_f\ \text{for every fine }p\}}$
is retained, where $t_f=t_{m+1}$. Thus its integrand is zero on the fine large-field
complement. The other sectors at the fine scale have already been attributed to the
large-field expansion.

Put $z=\widehat\beta_{m+1}$, $g=\beta'^{-1/2}$ and
$h=z^{-7/16}$. The letter $z$ is a scalar inverse coupling parameter in this
subsection; the Gaussian integration variable is $\xi$. Auxiliary
periods and coordinate counts are bounded by fixed powers of
$1+\log z$. Write $M=M_{m+1}^{\rm cut}$ for the cutoff amplitude
in \eqref{ch04:thresholds}, and take $D=d_0M$ with a fixed small
$d_0>0$. We shall use the following two distinct facts:
\[
 M=z^{o(1)},\qquad (1+\log z)^k=o(M^\epsilon)
 \quad\text{for every fixed }k,\epsilon>0.
\]
Both follow directly by taking logarithms of
$M=\exp((\log\log z)^2)$.

Here and below $z^{o(1)}$ in an upper bound means that for every
$\epsilon>0$ the quantity is $O_\epsilon(z^\epsilon)$ as $z\to\infty$
at fixed stage and block ratio.

Let $x$ be the coarse logarithmic coordinates, with their supremum
norm, and $y\in\RR^n$ the free fine coordinates after fixing the block axial
gauge. The classical action including the penalty term of the block kernel is
$\phi(x,y)$; write $\rho(x)$ for its critical configuration and
$A(x)=D_y^2\phi(x,\rho(x))$ for its positive Hessian.
The extra exponent $\psi(x,y)$ includes the fine loop term and minus
the logarithm of the product Haar-chart Jacobian. Suppose the critical-configuration map $\rho$,
action, extra exponent and relative logarithmic determinant have the
fixed derivative orders required below on a common neighborhood,
with bounds $z^{o(1)}$. The Hessian gap is uniform in the auxiliary
volume at the fixed block ratio. These are the stationary and path
estimates of Chapters~\ref{ch05:gaussian}--\ref{ch07:matching}; no
different flat covariance is substituted here.

For a fine-coordinate displacement $g\xi$, with $\xi\in\RR^n$, set
$s(x)=\tfrac12\log\det A(x)$ and define the second-order Taylor remainder
divided by $g^2$ through its integral formula
\[
 V(g,x,\xi)=\int_0^1(1-t)
 D_y^2\phi(x,\rho(x)+tg\xi)[\xi,\xi]\,dt.
\]
Stationarity gives
$\phi(x,\rho(x)+g\xi)-\phi(x,\rho(x))=g^2V(g,x,\xi)$,
including the limiting value
$V(0,x,\xi)=\tfrac12\xi^*A(x)\xi$.
With $\gamma_0=N(0,A(0)^{-1})$, put
\begin{align}
 a(g,x,\xi)&=s(x)-s(0)+V(0,0,\xi)-V(g,x,\xi)\notag\\
 &\quad+\psi(x,\rho(x))-\psi(x,\rho(x)+g\xi),\notag\\
 R(g,x)&=\int_{|\xi|\le D}e^{a(g,x,\xi)}\,d\gamma_0(\xi),
 \qquad L(g,x)=\log R(g,x).
 \label{ch09:truncated-integral}
\end{align}
The determinant factor is essential: at $g=0$ the exact identity is
\[
 R(0,x)=N(0,A(x)^{-1})\{\,|\xi|\le D\,\},
\]
which is the Gaussian measure of a ball, not identically one. Changing $\xi$ to
$-\xi$ also gives the exact parity $R(-g,x)=R(g,x)$.

Here are the uniform bounds needed for this elementary application
of Laplace's method. Differentiating the divided difference on
$|\xi|\le D$ costs a fixed polynomial in $D$ and the stationary
derivative bounds. It is therefore $z^{o(1)}$ at every fixed order.
The exponent vanishes at $(g,x)=(0,0)$, so on
$|g|,\|x\|_\infty\le h$ it is at most one in absolute value for
large $z$. If the reference Gaussian measure of the ball is at least $1/2$, then
$e^{-1}/2\le R\le e$, and the same derivative bounds $z^{o(1)}$ hold
for $L$. To obtain derivatives through order $r$ and the two
derivatives in $g$ used for parity, it suffices to retain derivatives
of $\phi,\rho$ through order $r+4$ and of the other terms through
order $r+2$. These are finite-stage requirements.

If $A(x)^{-1}\le cI$ and $n$ is at most a fixed polynomial in
$\log z$, a coordinate Gaussian tail bound gives
\[
 1-R(0,x)\le 2n\exp\bigl(-D^2/(2cn)\bigr).
\]
This is smaller than every fixed inverse power of $z$.
The same conclusion for its fixed derivatives follows either by
the Gaussian score formula or by finite-order interpolation on a
ball of radius $z^{-1}$ inside the parameter domain. We recall the
needed interpolation fact: a superpolynomially small value and
polynomially bounded derivatives through order $r+1$ imply
superpolynomially small derivatives through order $r$ on a smaller
ball. It follows successively from Taylor's estimate
$\|DF\|\le2\sup|F|/u+u\sup\|D^2F\|$ on a radius-$u$ ball,
applied to successive derivatives with shrinking radii.
No coordinate-dimension factor is needed when operator norms are
used. Parity and Taylor's theorem in $g$ now yield, for every fixed
$K>0$,
\begin{equation}
 \|D_x^r L(g,x)\|\le C_r(z)\beta'^{-1}+z^{-K},
 \qquad C_r(z)=z^{o(1)}.
 \label{ch09:even-gain}
\end{equation}
This includes $r=0$; the cutoff error at $g=0$ has been retained
throughout the argument.

\subsection{Identifying the retained region in the pure block integral}
\label{ch09:stationary-exterior}

Write $\mathcal D_x(y)$ for the configuration map of the block axial gauge.
Consider the region
$\mathcal S_x=\mathcal D_x(\rho(x)+g\overline B(0,D))$.
Let $\mathscr Z(x)$ be the pure block integral, and set
\[
 \Phi(x)=\beta'\phi(x,\rho(x))+\psi(x,\rho(x))
                         +\tfrac12\log\det A(x).
\]
If $N$ is the number of free fine links and $c_H$ the Haar density
constant, its normalization constant is
$\Omega=\pi_{\beta'}(\sqrt{2\pi}g)^n c_H^N$;
$\pi_{\beta'}$ is the block-kernel normalizing prefactor.
When $\mathcal S_x$ lies in the chart, all its fine plaquettes satisfy
the regularity condition $\chi(U_p)\le t_f$, and the actions for all values of the auxiliary index with positive probability weight coincide with $\phi$, the
change of variables gives exactly
\begin{equation}
 \int_{\mathcal S_x}\text{(pure integrand)}
                  =\Omega e^{-\Phi(x)}R(g,x).
 \label{ch09:retained-identity}
\end{equation}
The probability weights of the auxiliary index sum to one, so there is no additional factor counting its values.
The hypotheses about this region must be established before
using the identity.

Here is a verification using the geometry and the Hessian gap.
By the estimate \eqref{ch05:nonlinear-tree} in the block axial gauge, a fine
configuration lies in a fixed closed coordinate ball of the chart\mgubi{Issue 10: in which norm? existence/ takes the supremum ball and uses its convexity.}, the \emph{safe chart}, whenever its fine
plaquettes and the deviations of the coarse links from the values of the averaging operation are below a fixed $\delta>0$.
On that chart assume $D_y^2\phi\ge\gamma I/2$. Stationarity gives
\[
 \phi(x,y)-\phi(x,\rho(x))
                      \ge\gamma|y-\rho(x)|^2/4.
\]
Thus a chart point outside $\mathcal S_x$ adds at least
$\Delta=\gamma D^2/4$ to the exponent. A point outside the chart
with all fine plaquettes small must have such a deviation exceeding
$\delta$, and contributes a fixed positive multiple of $\beta'\delta^2$.
The possibly negative cubic, quartic and loop terms are bounded by
the path coefficient bounds. If $W$ bounds the auxiliary coordinate and site
counts, their total contribution to the exponent is at most
\[
 W B_{\rm loop},\qquad
 B_{\rm loop}=\beta'C(t_f^3+t_f^4)+\Lambda t_f^2,
 \qquad 0\le\Lambda\le2M.
\]
All such terms and the logarithmic Haar/determinant normalization
term are $o(M^2)$. In particular, the normalization term
$-N\log[c_H(\sqrt{2\pi}g)^{\dim\gLie}]$ is only a polynomial
in $\log z$. The bound on $\phi(x,\rho(x))$ on $\|x\|<h$
also leaves the fixed penalty for a deviation exceeding $\delta$ available, because
that value tends to zero and $\Delta/\beta'\to0$.
Consequently the exterior integrand is bounded by
$\Omega e^{-\Phi(x)}e^{-\Delta+o(M^2)}$.

It remains to verify the regularity condition on the fine plaquettes on the retained ball. For the
auxiliary configurations used here, which are reflected periodic extensions (pullbacks of the original field to an auxiliary torus under the coordinatewise reflection map $\pi_v$ of Section~\ref{ch09:local-density}), this can be done directly
from the induced map on curvatures, without a lower bound from a sector
expansion. A full coordinate cycle of a reflected periodic extension traverses an
original path and then its inverse. Its holonomy is therefore exactly
one at every transverse base point. Equivalently, choose a coordinate
axial gauge on the original cube and pull that gauge back under the
reflection map. Plaquette filling bounds give logarithmic coordinates
with
\[
 \|x\|_\infty\le P(B)t_c,\qquad
 \|d_{1,c}x\|_\infty\le C t_c+P(B)t_c^2,
 \qquad t_c=t_m.
\]
The second estimate follows by expanding each plaquette product in
these small link coordinates. Crucially, its leading constant is
independent of $B$; the polynomial loss occurs only in the quadratic
error. This argument would fail for an arbitrary small-curvature
periodic field with nontrivial winding holonomies.

Let $\mathcal U(x)=\mathcal D_x(\rho(x))$ be the critical configuration as a physical
fine field. Its linearization is the representative $W_F$ in the block axial gauge of
Chapter~\ref{ch05:gaussian}. Its curl agrees with that of the smooth
representative $W_a$. The curvature identity and the
fine-row bound of the induced map $T$ on curvatures give
\[
 d_{1,f}W_a=H_f^{-1}TH_c d_{1,c},\qquad
 \|H_f^{-1}TH_c\|_{\infty\to\infty}\le C_jL^{-2}.
\]
The operator norm here sums each fine-output row over coarse
inputs; the opposite column sum has a different scaling.
The quadratic remainder of the critical-configuration map and the plaquette
product is bounded using their common second derivatives. Thus
\begin{equation}
 \sup_p\chi(\mathcal U(x)_p)
       \le C_jL^{-2}t_c+P_{j,L}(B)t_c^2.
 \label{ch09:stationary-cutoff-margin}
\end{equation}
Choose $L$ so that $C_jL^{-2}t_c\le t_f/8$, using the uniform
bound $t_c/t_f\le2$. Then increase the lower bound on $z$ so that
the second term is at most $t_f/8$. These choices are in the
required order: the constant in the first term is fixed before
$L$, while the polynomial in the second term may depend on $L$.

Plaquette functions have a fixed Lipschitz constant $L_0$ in the
physical fine-link supremum coordinates. Local chart maps have
fixed-stage derivative bounds. Consequently variation of $x$ by
$z^{-1}$ and of the fine coordinates at the critical configuration by $g\xi$,
$|\xi|\le D$, changes every plaquette distance from the identity by at most
$C_{j,L}z^{-1}+L_0gD$, after incorporating the fixed chart norm
conversion into $L_0$. Choose $d_0$ sufficiently small, so that
$L_0gD\le t_f/8$, and then raise the lower bound on $z$ for the term
$C_{j,L}z^{-1}$. Every point of each retained region now satisfies the regularity condition on the fine plaquettes,
with a strict margin. The common chart and the ball on which the actions for all values of the auxiliary index with positive probability weight coincide
contain these regions because $h+gD\to0$.
We have therefore obtained, on a full microscopic neighborhood of
each tested coarse configuration satisfying the regularity condition,
\begin{equation}
 0\le\mathscr Z(y)-\Omega e^{-\Phi(y)}R(g,y)
       \le\Omega e^{-\Phi(y)}e^{-c_0M^2},\qquad c_0>0.
 \label{ch09:exterior-bound}
\end{equation}
This deduction uses the linearization of the critical configuration (through its smooth representative $W_a$), its second
derivatives, the compatible path bounds and the geometric
estimate \eqref{ch05:nonlinear-tree}. The direct proof of the regularity condition applies to the
auxiliary fields that are reflected periodic extensions and to their microscopic neighborhoods, which
are precisely the configurations used for localization.

\subsection{Centering the full logarithm}
\label{ch09:centered-log}

Put $E=\log\mathscr Z-\log\Omega+\Phi-L(g,\cdot)$.
Since $R\ge e^{-1}/2$, \eqref{ch09:exterior-bound} gives
$0\le E\le2e\,e^{-c_0M^2}$. To pass to derivatives, write
the positive pure integrand at a fixed fine field and a fixed value of the auxiliary index as
$a(\omega)e^{p(x,\omega)}$, with $a\ge0$ independent of $x$.
Only the coarse-link kernel is differentiated. Its exponent has
fixed-order derivatives bounded by $C_r\beta' n_b$, where
$n_b$ is the number of coarse links. The exponential chain rule
and integration give relative bounds
\[
 \|D^r\mathscr Z(x)\|\le P_r(\beta'n_b)\mathscr Z(x).
\]
The logarithm chain rule thus gives polynomial bounds without
dividing an absolute error by a possibly tiny partition function.
The other terms in $E$ have polynomial derivative bounds from the
stationary estimates. Interpolation on the microscopic neighborhoods
in \eqref{ch09:exterior-bound} proves that every fixed jet of $E$
is smaller than every inverse power of $z$.

Let $a_c=\mathcal K$ be the coarse classical action and
$d_c=\mathcal L+\alpha\mathcal Q$ the prescribed coarse quantum contribution:
the inherited loop action of the coarse path action and its quadratic part
multiplied by the coefficient $\alpha$ fixed by the matching, as in
\eqref{ch09:early}.\mgubi{Identified from the Lean definition YM.\allowbreak Precision.\allowbreak draft ($\alpha$ times the quadratic term plus the loop term), existence/\allowbreak stationary-output-construction.\allowbreak tex and v2 D-section; no $\Delta$ and no constant per site. Open: Lean uses $\mathcal K$, $\mathcal L$, $\mathcal Q$ without event indicators, while v3 defines them with indicators; they coincide on the small-field ball.} Define
\[
 C=a_c-\phi(\cdot,\rho),\qquad
 W=d_c-\psi(\cdot,\rho)-\tfrac12\log\det A.
\]
The corrected logarithm has the exact decomposition
\begin{equation}
 \log\mathscr Z+\beta'a_c+d_c
                  =\log\Omega+L+E+\beta'C+W.
 \label{ch09:pure-decomposition}
\end{equation}
Classical matching gives $D^rC(0)=0$ for $0\le r\le4$.
The quadratic quantum matching bounds $D^2W(0)$ by an
exponentially small finite-period error, with polynomial coordinate
norm costs. Its higher fixed derivatives are $z^{o(1)}$.
The full corrected logarithm has zero first derivative at zero
by constant gauge invariance and simplicity of $\gLie$.
Thus $DW(0)=-DL(0)-DE(0)$; no separate parity claim about each
determinant contribution is needed.

For a function on the radius-$h$ ball write its scaled derivative
bound through order $r$ as
$\max_{0\le k\le r}h^k\sup\|D^kf\|$.
Taylor's theorem, \eqref{ch09:even-gain}, and the two matching
cancellations show that the left side of
\eqref{ch09:pure-decomposition} minus its value at zero has this bound at most
\begin{equation}
 C_r(z)\bigl(\beta'^{-1}h+\beta'h^5+h^3\bigr)+z^{-K},
 \qquad C_r(z)=z^{o(1)},
 \label{ch09:pure-small-bound}
\end{equation}
where $K$ is arbitrary and the exponentially small quadratic
finite-period error has been included in the last term.
For $h=z^{-7/16}$ the three powers are respectively
$z^{-23/16}$, $z^{-19/16}$ and $z^{-21/16}$.
Their largest exponent is $-19/16<-23/20$.
Original small-field logarithmic coordinates and the directions admitted by the activity
seminorm \eqref{ch04:local-seminorm}, restricted to the cube $N_{\bar w}$ about the localization site and carried to the reflected periodic extension ($\bar w$ and the reflected extension are those of Section~\ref{ch09:local-density} below; the cube $N_r(x)$ is used there without a definition), are bounded by a polynomial in $B$ times
$t_m$, which is $o(h)$.\mgubi{All directions of the seminorm are meant, not only the affine ones (v2, D-section). v2 adds that the reflected torus must admit the charts and derivative bounds used here, and says that this and the finite-period comparison on a common analytic strip ``remain under verification''; v3 does not prove it either.} They therefore obey the same estimate.
This proves the auxiliary remainder needed for localization from
the explicit stationary, matching and actual-sector inputs above.
It does not require a formal expansion of an untruncated
exponential of a cubic Gaussian polynomial.


\subsection{Recovering a local density}

The passage from auxiliary sums to local activities has an
exact finite-dimensional explanation. Truncate the selected
low-order formula, retaining a weighted tail
$\varepsilon_{\rm tr}$. Give an ordinary selected term $f_a^0$
the shares $p_{xa}=\mathbf1_{X_a}(x)/|X_a|$, and a path term
the share $\mathbf1_{\{o_a\}}(x)$. These shares sum to one.
Set
\[
 F_x=\sum_{a\ {\rm selected}}p_{xa}f_a^0,\qquad
 M_x=\sum_{a\ {\rm selected}}p_{xa}h_{X_a}f_a^0.
\]
Choose a common dependence radius $R$ about every site receiving
a nonzero share. Set $w=4R+100(1+B)$, $\bar w=w+5$, and use
auxiliary tori $\Lambda_i$ of sides $2(w+i)$, $0\le i\le4$.
The relevant reflection cubes must embed in the original torus.
For a cyclic coordinate $q$ of period $2v$, define
$\pi_v(q)=\min(\bar q,2v-\bar q)$, where $0\le\bar q<2v$.
Apply this map coordinatewise. Write $U_{i,z}$ for the pullback of
the original field under $q\mapsto z+\pi_{w+i}(q)$ on $\Lambda_i$.
Reversed links are inverted, so small plaquette distances from the identity stay small.
Let $F_y^{(i)}$ denote the same local formula on $\Lambda_i$
with its flat value subtracted, and set
\begin{align}
 \lambda_i&=\frac{(-1)^i\binom4i}{2^4\,4!},&
 G_x&=\sum_{i=0}^4\lambda_i\sum_{y\in\Lambda_i}
                              F_y^{(i)}(U_{i,x}),\notag\\
 L_x&=M_x-F_x(1)-
   \sum_{i=0}^4\lambda_i\sum_{y\in\Lambda_i}
       h_{N_{\bar w}(x-\pi_{w+i}(y))}
                      F_y^{(i)}(U_{i,x-\pi_{w+i}(y)}).
 \label{ch09:localization}
\end{align}
In the second line the original cube is based at
$x-\pi_{w+i}(y)$, so its site
$y$ corresponds to the original site $x$.
This dependence on $y$ is essential.

When $N_{2\bar w}(x)$ is small-field, all these event indicators are one.
Classify sites of $\Lambda_i$ by their distances, through $R$,
from the reflection faces. With the recentering just specified,
the local expression in each class is independent of $i$.
The multiplicity of each class is polynomial in $i$ of
degree at most four; only the all-interior class has degree
four, with leading coefficient $2^4$.
The fourth-difference weights therefore give
$\sum_i\lambda_i\sum_y F_y^{(i)}(U_{i,x-\pi_{w+i}(y)})
=F_x(U)-F_x(1)$, and $L_x=0$.
This is not a pointwise assertion that $G_x=F_x-F_x(1)$.

Sum \eqref{ch09:localization} over $x$ and change variables
$x-\pi_{w+i}(y)\mapsto x$ in each finite auxiliary sum.
Translation covariance makes $F_x(1)=F_0(1)$, giving exactly
\[
 \sum_x M_x=|\Lambda|F_0(1)+
              \sum_x h_{N_{\bar w}(x)}G_x+\sum_xL_x .
\]
Thus $G_x$ is the small-field local activity and $L_x$ is a
logarithmic correction supported where its enlarged
neighborhood contains a large-field site.

Auxiliary volumes and reflected derivative norms grow only
polynomially in $B$. Thus \eqref{ch09:aux-input} implies
\[
 \|G_x\|_m\le P_j(\mathfrak b)
           (\mathfrak b^{-23/20}+\varepsilon_{\rm tr}),
 \qquad \sup|L_x|\le P_j(\mathfrak b)\mathfrak b^{-7/16}.
\]
The first fits an arbitrarily small fraction of
$\epsilon_m^{\rm reg}$, since $23/20>11/10$.
The second tends to zero and is incorporated into
large-field activities by the exact operation below.

The selected low-order terms contain no incoming regular or
large-field activity. Their coefficients depend only on the
prescribed path, actual inverse coupling and fixed labels. Their
weighted continuity on the compact inverse coupling band therefore
permits one truncation for all inputs with those data.
Terms involving incoming activities have already been
estimated separately. More generally compactness along a
single activity curve would give a truncation only along
that curve. All choices here are retained under restriction
of the parameter interval.

\section{A precise preparatory output}
\label{ch09:preparatory}

The preceding estimates, when their concrete local inputs
hold, give local functions $f_i$, preliminary large-field
activities $k_i$, and a scalar $\tau$ satisfying
\begin{equation}
 \mathscr Z_{\rm blk}(U)=
 e^{|\Lambda|\tau-\beta'\mathcal K(U)-\mathcal L(U)
                       -\alpha\mathcal Q(U)-\sum_i h_{X_i}f_i(U)}
                         \mathcal F^0(U).
 \label{ch09:early}
\end{equation}
Here $\mathcal F^0$ is the exact-cover sum of the $k_i$,
and $\mathcal Q$ is the quadratic part, restricted by its event indicators, of the
prescribed path action. The coefficient $\alpha$ is
the matched determinant coefficient, not an adjustable
trial inverse coupling. The bounds needed below are
\begin{equation}
 \sup_x\sum_{o_i=x}e^{(\sigma+3)c(X_i)}a_i\le H_{\rm pre},
 \quad a_i\ge\|f_i\|_m,\qquad
 \sup_x\sum_{i:x\in Y_i}e^{(\sigma+1)c(Y_i)}k_i^{\rm maj}
       \le e^{-c_j(M_{m+1}^{\rm cut})^2},
 \label{ch09:early-bounds}
\end{equation}
where $c(X)=\operatorname{cov}_{B_m}(X)$ and
$H_{\rm pre}=\theta c_j^{\rm pre}\epsilon_m^{\rm reg}$,
$0<\theta\le1$. The positive stage constant $c_j^{\rm pre}$
is fixed from the subsequent subtraction and path-replacement
operator norms before choosing the ratio.

To obtain these bounds choose $L$ first so that
$V_jL^{-1}\epsilon_{m+1}^{\rm reg}$ fits the desired fraction of
$H_{\rm pre}$. This is a condition on the fixed stage constants:
the ratio $\epsilon_{m+1}^{\rm reg}/\epsilon_m^{\rm reg}$ is bounded
on the permitted neighboring reference-inverse-coupling bands.
Then choose the inverse coupling threshold for the inherited
comparison, \eqref{ch09:degree-gain},
\eqref{ch09:aux-input}, and the small logarithmic corrections.
Choose the truncation accuracy after its polynomial
localization costs have been specified.
Allocate distinct positive fractions of the exponent in
\eqref{ch09:joint-input} to sector multiplicities, local
normalization losses, and the final large-field bound.
Polynomial support costs use support-weight reserves.
An unexplained extensive normalization cannot be absorbed into
a single selected large-field condition.

Equation \eqref{ch09:early} and its bounds are the concrete
interface for the rest of the chapter. The remaining
operations and estimates will now be proved directly from it.

\section{Local quadratic subtraction}
\label{ch09:subtraction}

For $a<b$ define the clover curvature at $o$ by
\[
 F^{\rm cl}_{ab}(U,o)=\tfrac14\sum_{\varepsilon,\eta=\pm1}
 \varepsilon\eta\,\operatorname{proj}_{\gLie}
       \frac{U_{o;\varepsilon a,\eta b}
                         -U_{o;\varepsilon a,\eta b}^{-1}}2 .
\]
The loop follows the four signed edges and starts at $o$.
This field transforms by $\Ad$ at $o$ and reads a radius-two
neighborhood. For an affine potential
$A_\mu(y)=u_\mu+\sum_\nu y^\nu v_{\mu\nu}$ its derivative
at the identity is $v_{ba}-v_{ab}$.

For a preliminary $f$ on $X$ anchored at $o$, put
$Z=X\cup N_2(o)$. On a nonwrapping $Z$, let
$A^{q,T}_c(y)=\mathbf1_{\{c=b\}}\xi_o(y)^aT$ for $q=(a,b)$,
and fix a unit $T_*\in\gLie$. Define
\[
 q_f(q,r)=D^2[f(\exp\cdot)](0)[A^{q,T_*},A^{r,T_*}],
 \quad
 Q_f(U)=\tfrac12\sum_{q,r}q_f(q,r)
       \langle F^{\rm cl}_q(U,o),F^{\rm cl}_r(U,o)\rangle.
\]
Set $q_f=0$ if $Z$ wraps. Gauge invariance makes the first
derivative at the identity vanish and the Hessian annihilate
local gradients. The zero-curl part of an affine potential
is a discrete gradient of a quadratic site function.
It can be extended periodically using a cutoff outside
the dependence region. Therefore the Hessian on affine fields
depends only on their constant curls. Invariance under the adjoint action and simplicity of $\gLie$ make its
bilinear form a multiple of the invariant metric.
Thus $Q_f$ has the same affine Hessian as $f$.
This is the local Ward-identity argument of
Section~\ref{ch07:kernels} in finite coordinates.

Set
\[
 R_f=f-f(1)-Q_f,\qquad
 \zeta_f=(1-h_Z)\{f(1)+h_{\{o\}}Q_f-h_Xf\}.
\]
An exact two-case identity is
\begin{equation}
 h_Xf=f(1)+h_{\{o\}}Q_f-\zeta_f+h_ZR_f .
 \label{ch09:subtract-identity}
\end{equation}
On a small-field $Z$ all event indicators are one; otherwise $h_Z=0$
and the definition of $\zeta_f$ proves the equality.
The remainder vanishes at the flat field and has zero
affine Hessian on nonwrapping supports.
The function $\zeta_f$ vanishes on small-field $Z$.
Its large-field indicator dependence must still be recorded in its full
dependence region; no derivative of an event indicator is taken.

The evaluations in normalized affine directions and polarization give
$t^2|q_f(q,r)|\le3a_f$. Every allowed derivative of
$F^{\rm cl}$ is $O_j(t)$ on its small-field domain, so
\begin{equation}
 \|R_f\|_m+\sup|\zeta_f|\le C_j a_f .
 \label{ch09:subtraction-bound}
\end{equation}
In the supremum term, $f$ and $Q_f$ occur only under their
own small-field conditions. Enlarging $X$ to $Z$ costs a bounded
number of covering balls.

The conversion of anchored bounds to containing-site bounds
is explicit. If $x\in Z$ and $c(Z)=n$, its anchor lies
within distance $CBn$ of $x$. There are at most
$C(1+B)^4(1+n)^4$ choices. An additional exponent in the support weight
absorbs this polynomial. Consequently the discrepancy
series has containing-site norm at exponent $\sigma+1$
bounded by $C_j(1+B)^{d_j}H_{\rm pre}$.

By translation and hypercubic covariance,
\[
 \sum_{o_i=o}q_{f_i}(q,r)=\Delta\delta_{qr}.
\]
An axis reflection kills off-diagonal orientation pairs;
permutations identify the diagonal pairs.
Absolute convergence follows from the derivative bound,
and
\begin{equation}
 t^2|\Delta|\le C_jH_{\rm pre}.
 \label{ch09:delta-bound}
\end{equation}
Thus, with
$\mathcal C(U)=\frac12\sum_o h_{\{o\}}\sum_q|F_q^{\rm cl}(U,o)|^2$,
the sum of the extracted quadratic functions is
$\Delta\mathcal C(U)$.

\section{Replacing the local quadratic term by the path action}
\label{ch09:path-reset}

The path action is the tested and clipped density of \eqref{ch04:path-action}. Its path and cutoff data remain fixed during this replacement, as do the frame choices defining the actual probability kernel \eqref{ch04:block-operator}.

Expand the quadratic part into plaquette channels and ordered
pairs of entries in each row square. For a channel $d$ let
$a_dP_d$ be its quadratic expression without its event indicator, $a_dQ_d$
its constant-curvature clover substitute, and $L_d$ its
actual clipped or tested expression. Its support $Y_d$
contains both paths and every nonautomatic relaxed small-field condition.
On the small-field set of $Y_d$, $L_d=a_dP_d$.
Define
\[
 F_d=\Delta a_d(Q_d-P_d),\qquad
 J_d=(1-h_{Y_d})\{(\alpha+\Delta)L_d
                              -\Delta a_dh_{\{o_d\}}Q_d\}.
\]
Then
\begin{equation}
 (\alpha+\Delta)L_d+h_{Y_d}F_d
 =\alpha h_{Y_d}a_dP_d+
       \Delta a_dh_{\{o_d\}}Q_d+J_d .
 \label{ch09:channel}
\end{equation}
The equality follows on each of the two health domains.
Every $F_d$ has zero affine Hessian and every $J_d$
vanishes on its small-field support.

The selected free metric satisfies
$c_{\rm dir}I+S(0)^*S(0)=I$, where $S(0)$ is the row
coefficient matrix at zero momentum.
Therefore summing the clover substitutes gives exactly
$\mathcal C$, while the raw channels restricted by their event indicators sum to
$\mathcal Q$. Add the regular terms
$-(\alpha+\Delta)A_d$ for the cubic and quartic atoms.
After the common signed-symmetry average, let $F$ be the
resulting sum of regular terms restricted by their event indicators and $J$ the sum of the $J_d$.
Equation \eqref{ch09:channel} becomes
\begin{equation}
       (\alpha+\Delta)\mathcal K+F
                   =\alpha\mathcal Q+\Delta\mathcal C+J.
 \label{ch09:path-identity}
\end{equation}
Combining \eqref{ch09:early},
\eqref{ch09:subtract-identity} and
\eqref{ch09:path-identity} gives, at every field,
\begin{equation}
 \mathscr Z_{\rm blk}=
 e^{|\Lambda|\tau-\sum_i f_i(1)
       -(\beta'+\alpha+\Delta)\mathcal K-\mathcal L-E}
       \mathcal F^0\,e^{\sum_i\zeta_{f_i}+J},
 \qquad E=\sum_i h_{Z_i}R_{f_i}+F .
 \label{ch09:before-fusion}
\end{equation}
No division by the exact-cover gas was used, so this equality
also holds where that gas vanishes or is signed.

Weighted path sums and additional exponential decay in path length give constants
$C_{\rm reset},L_0,L_1,L_h$ such that
\begin{align}
 \mathcal E_m(E)&\le C_{\rm reset}H_{\rm pre}
     +C_j\{H_{\rm pre}L_0+
                  (|\alpha|t^3+H_{\rm pre})L_1\},\\
 \|\zeta+J\|_{\sigma+1,{\rm site}}
 &\le C_j(1+B)^{d_j}H_{\rm pre}
   +C_j\{t|\alpha|+H_{\rm pre}/t+H_{\rm pre}\}L_h .
 \label{ch09:reset-cost}
\end{align}
Here the site norm sums supremum majorants of logarithmic
terms over supports containing a site.
The low regular channels contain two curvatures, the high
channels at least three; each derivative costs a polynomial
in path length, controlled by the exponential path norm.
The actual clipped channel can have size $O(t)$ on a
large-field configuration, which explains the term
$t|\alpha|$ in the second line.

Choose $c_j^{\rm pre}$ in \eqref{ch09:early-bounds}
small enough to absorb the terms linear in $H_{\rm pre}$
in the first line. At fixed ratio $|\alpha|$ is bounded,
and
\[
 t^3/\epsilon_m^{\rm reg}
       =\widehat\beta_m^{-2/5}(M_m^{\rm cut})^3\longrightarrow0.
\]
Thus the first line is at most $\epsilon_m^{\rm reg}$.
The second line tends to zero even after multiplication
by $(2B+1)^4$: in particular
$H_{\rm pre}/t=O(\widehat\beta_m^{-3/5}/M_m^{\rm cut})$.
This is the smallness needed for the exact regrouping.

\section{Absorbing logarithmic corrections into large-field activities}
\label{ch09:fusion}

Write the logarithmic corrections in
\eqref{ch09:before-fusion} as $\ell_a$ on nonempty exclusion
supports $Z_a$, with $\ell_a=0$ when $Z_a$ is small-field.
Their extra large-field indicator dependence regions are retained.
For each finite connected selection $P$ of preliminary
large-field labels and logarithmic labels, require a
nonempty large-field part and pairwise disjoint preliminary
exclusion supports. Connectivity refers to intersections
of these exclusion supports and the $Z_a$.
Set $Y_P$ equal to their union and $I_P$ equal to the union
of the preliminary prescribed large-field site sets. Define
\[
 K_P(U)=\mathbf1_{\{Y_P\cap\operatorname{Bad}_m(U)=I_P\}}
     \prod_{i\in P_{\rm lf}}k_i(U)
     \prod_{a\in P_{\log}}(e^{\ell_a(U)}-1),
\]
with the conjunction of all inherited additional conditions.

\begin{proposition}[Exact regrouping]
\label{ch09:fusion-prop}
Suppose $\sum_a\sup|\ell_a|<\infty$. The exact-cover sum of
the $K_P$ equals $\mathcal F^0 e^{\sum_a\ell_a}$.
If nonnegative majorants satisfy
\[
 |k_i|\le e^{-\kappa}v_i,\qquad
 |e^{\ell_a}-1|\le v_a,\qquad
 \sup_x\sum_{a:x\in Z_a}e^{(\sigma+1)c(Z_a)}v_a\le\gamma,
 \quad (2B+1)^4\gamma\le1,
\]
where the unified index also includes the preliminary
large-field supports, then
\[
               \mathcal K_m(K)\le e^{-\kappa}(e^\gamma-1).
\]
\end{proposition}

\begin{proof}
Expand $e^{\sum_a\ell_a}=\prod_a[1+(e^{\ell_a}-1)]$.
In a nonzero term each selected logarithmic label contains
a large-field site. Exact coverage assigns that site to a preliminary
prescribed large-field site set, so its connected component contains
a large-field label. Taking components gives disjoint union
supports and the displayed new prescribed-site equality.
Conversely flattening a compatible family of new labels
recovers the old exact cover and the logarithmic selection.
No constituent label is repeated because its support is
nonempty. All additional conditions are retained. This is a bijection
of finite terms and does not require independence of their
large-field indicator reads.

For the bound retain $e^{-\kappa}$ from one large-field
label per component, and bound the others by $v_i$.
Subadditivity of $c$ on connected unions and the rooted
tree bound underlying Theorem~\ref{ch08:kp} sum the
connected selections through a fixed site. One additional unit of exponential support decay bounds the number of attachments, since a support has at most
$(2B+1)^4c(Z)$ sites. Their sum is bounded by
$\gamma$, hence by $e^\gamma-1$.
The resulting absolute envelope justifies the countable
regrouping.
\end{proof}

Take $\kappa$ to be half of the exponent retained in the
second bound of \eqref{ch09:early-bounds}.
The remaining half makes the preliminary large-field
majorants small in the attachment norm.
For the logarithmic terms use
$|e^\ell-1|\le e^{\sup|\ell|}\sup|\ell|$ and
\eqref{ch09:reset-cost}.
Then $\gamma\to0$ and a strict margin
$(2B+1)^4\gamma\le1/2$ is available.
Finally
\[
 e^{-\kappa}(e^\gamma-1)
       \le \exp(-(M_m^{\rm cut})^{1/5})
             =\epsilon_m^{\rm lf}
\]
for a sufficiently large inverse coupling threshold, provided the
neighboring reference scales have the prescribed comparison.
The square $(M_{m+1}^{\rm cut})^2$ in $\kappa$ dominates
the exponent $1/5$ and all polynomial logarithmic costs.

\section{The normalization of the field-independent term and continuity}
\label{ch09:continuity}

The output inverse coupling and scalar are now fixed:
\begin{equation}
 \beta=\beta'+\alpha+\Delta,\qquad
 v_{\rm sc}=\tau-|\Lambda|^{-1}\sum_i f_i(1).
 \label{ch09:reset}
\end{equation}
Equation \eqref{ch09:before-fusion} and
Proposition~\ref{ch09:fusion-prop} prove
$\operatorname{dens}(\rho)=\mathscr Z_{\rm blk}$ pointwise.
Flat normalization also gives
$v_{\rm sc}=|\Lambda|^{-1}\log\mathscr Z_{\rm blk}(1)$.
The exact block integral of a positive input is positive;
signed individual activities cause no contradiction.

Projection already provides the useful estimate
\[
 |\Delta|\le C_jH_{\rm pre}/t_m^2
       \le C_j\theta c_j^{\rm pre}
                   \widehat\beta_m^{-1/10}(M_m^{\rm cut})^{-2}
       \le1 .
\]
If the matched coefficient has
$\alpha=-\kappa_G\log L+O_j(1)$ as in
Proposition~\ref{ch07:slope}, then
\begin{equation}
       \beta=\beta'-\kappa_G\log L+e,\qquad |e|\le d_j+1.
 \label{ch09:drift}
\end{equation}
No estimate of $\Delta$ was used to define its local
projection or prove the exact algebra.

There is a second, intrinsic way to identify the reset.
On an even-sided torus choose direction-one links alternately
equal to $u$ and $1$ on even and odd second-coordinate
slices, with $\chi(u)=rt$ for sufficiently small fixed $r>0$.
Every nontrivial plaquette has this distance from the identity. The direct
positive plaquette term dominates the cubic and quartic
terms when $r$ is small, so this small-field configuration
$U_*$ satisfies $\mathcal K(U_*)\ge c_*|\Lambda|t^2$.
Hold the constructed activities fixed and vary only a trial
inverse coupling correction $d$. The scalar discrepancy is
\[
 f(d)=(\beta'+\alpha+d)\mathcal K(U_*)+\mathcal L(U_*)+
 E(U_*)-\log\frac{\mathscr Z_{\rm blk}(1)}{\mathscr Z_{\rm blk}(U_*)}
       =(d-\Delta)\mathcal K(U_*).
\]
It has the unique zero $d=\Delta$.
Two evaluations identify the two scalars only because the
full-field identity has already been established.
Matching two fields alone would not prove that identity.

For a compact continuous input curve choose fixed local
supports, cutoff conditions, truncations and reflected tori.
Projection, subtraction and channel replacement are bounded
linear operations at these fixed labels. Products are
continuous by telescoping with one difference factor.
For changes of inverse coupling hold physical field coordinates
fixed; for a centered Gaussian of precision $\beta A$,
\[
 \frac{d\nu_{\beta_2}}{d\nu_{\beta_1}}(a)
  =(\beta_2/\beta_1)^{D/2}
             e^{-(\beta_2-\beta_1)\langle a,Aa\rangle/2}.
\]
Its derivative is the same density ratio times
$D/(2\beta_2)-\langle a,Aa\rangle/2$.
On a small relative interval it is dominated by a
polynomial times a Gaussian with a smaller positive
precision. Mean and matrix changes give analogous
Gaussian scores. Apply the common-law estimate to these
scores; do not differentiate a merely measurable incoming
large-field activity.

Infinite regrouping requires more than coefficientwise
continuity. Given a sequence of parameters approaching
$u_0$, pass to a subsequence whose weighted input
discrepancies have sum at most $\varepsilon$.
Their sum added to the majorant at $u_0$ is a common
summable envelope. The strict attachment margin allows
this addition for small $\varepsilon$.
The rooted estimate then gives one summable output envelope;
dominated convergence proves weighted continuity on this
subsequence. Every potentially nonconvergent sequence
would have such a convergent further subsequence, a
contradiction.

At fixed volume, the exact-cover sums obey the elementary
telescoping estimate
\[
 |\mathcal F_u(U)-\mathcal F_{u_0}(U)|
 \le e^{\sum_i M_i}\sum_i\sup|K_{i,u}-K_{i,u_0}|,
 \quad M_i=\max(\sup|K_{i,u}|,\sup|K_{i,u_0}|).
\]
Thus the densities are continuous uniformly in $U$.
The positive block-kernel normalizers are continuous and
bounded away from zero on compact inverse coupling intervals;
dominated convergence in probability Haar measure proves
continuity of $\mathscr Z_{{\rm blk},u}$.
The scalar formula and projection show continuity of
$v_{\rm sc}$ and $\beta$. Restriction to a closed
subinterval retains the construction and all its labels.

\section{Wilson initialization}
\label{ch09:wilson}

The finest-scale representation can be constructed directly
from the Wilson action \cite{Wilson1974}, without a limiting
argument. Write $W(u)=1-\operatorname{Re}\tr(u)/n_G$ for the Wilson
plaquette energy and $F(u)=\operatorname{proj}_{\gLie}(u-u^{-1})/2$.
The map $X\mapsto F(e^X)$ has derivative $I$ at zero and a
local inverse $\psi$. The map and its inverse are odd, and
$W(u^{-1})=W(u)$, so $W(e^{\psi(Y)})$ is even.
Let $T_4$ be the degree-four Taylor
term of the Wilson plaquette energy evaluated at
$e^{\psi(Y)}$. Then
\[
 R(u)=W(u)-\tfrac12|F(u)|^2-T_4(F(u))
\]
vanishes through degree four at the identity.
Choose the finest row
$\sqrt{1-c_{\rm dir}}F$, no cubic or loop term, and quartic
atom $(1-c_{\rm dir})T_4$ with its prescribed event indicator.

For a plaquette $p$ write $h_p$ for the health event indicator on its
vertices, $\chi_p^{\rm rel}$ for its row cutoff condition, and
$k_p=(1-c_{\rm dir})|F(U_p)|^2/2$.
Define the residual energy
\[
 \mathcal W_p=
 \begin{cases}
 W(U_p)-c_{\rm lin}t\chi(U_p),&\chi(U_p)>t,\\
 h_pk_p+(1-h_p)(1-c_{\rm dir})W(U_p)
                   -\chi_p^{\rm rel}k_p,&\chi(U_p)\le t.
 \end{cases}
\]
Since $|F(u)|\le\chi(u)$ and $W(u)=\chi(u)^2/2$,
this is nonnegative; when $\chi(U_p)>t$ it is at least
$t^2/4$. The relaxed small-field condition implies that the plaquette
itself is small, and the condition can also fail because of a nearby bad
plaquette. Direct substitution
gives
\[
 W(U_p)=L_p+(1-c_{\rm dir})h_pT_4(F(U_p))
           +(1-c_{\rm dir})h_pR(U_p)+\mathcal W_p ,
\]
where $L_p$ is the actual clipped quadratic channel.
If the relaxed small-field condition inspects a larger neighborhood, its
failure is included in $\mathcal W_p$ by the same formula;
the identity remains valid because on a small-field plaquette
the residual is $k_p-\chi_p^{\rm rel}k_p$.

Use $\beta(1-c_{\rm dir})R(U_p)$ as a regular activity.
Its normalized derivatives are bounded by $C_j\beta t^5$.
For orders at most five, the missing powers after
differentiation are supplied by Taylor's formula; for
higher fixed orders, the direction factors supply at
least $t^5$. Hence
\[
 \beta t^5/\epsilon_N^{\rm reg}
              =O(\widehat\beta_N^{-2/5}(M_N^{\rm cut})^5)\to0 .
\]
Assign the nonnegative residual energies to connected
components of large-field sites whose radius-$2B$ neighborhoods
overlap. A component site set $I$ has exclusion support
$Y_I=N_{2B}(I)$ and activity
\[
 K_I=\mathbf1_{\{Y_I\cap\operatorname{Bad}_N=I\}}
                         e^{-\beta d_I},
\]
where $d_I$ is the sum of assigned residual energies with
the signed average. Each large-field site meets a large-field plaquette,
and a plaquette has at most four vertices, so
$d_I\ge c\,t^2|I|$ for a fixed $c>0$.
Every residual term is assigned exactly once; its relaxed small-field condition has a witnessing large-field site within distance $B$.
The only nonzero exact-cover selection is the actual set
of these components, so the density identity is exact.

A connected prescribed set of $r$ large-field sites is bounded in number
by a rooted spanning-walk count. The containing-site
weighted sum is at most
\[
 C_j(1+B)^4\sum_{r\ge1}
       \left[C_j(1+B)^8e^{5^4\sigma-c\beta t^2}\right]^r .
\]
Since $\beta t^2\ge(M_N^{\rm cut})^2/2$, this fits
$\epsilon_N^{\rm lf}$ for a sufficiently large threshold.
The activity labels and eligibility criteria are independent of $\beta$.
The resulting full Wilson curve is continuous, and its
scalar normalization is zero.

\section{The one-step conclusion and its input data}
\label{ch09:step}

\begin{theorem}[One renormalization step]
\label{ch09:step-theorem}
Assume the derivative orders and spatial exponents stated at
the start of the chapter. Fix an adjacent quadratic pair with the
positive metrics, exact transport identities and actual gauge-covariant block average
of Chapter~\ref{ch05:gaussian}. Let the incoming path have the
normalized quadratic coefficient and exponential coefficient and
derivative bounds of Chapters~\ref{ch04:blocking} and
\ref{ch07:matching}. The incoming activities obey the small-field
cancellations, large-field eligibility rules and weighted bounds of
Chapter~\ref{ch04:blocking}, with the support reserve used above.
Assume also that the represented input density is strictly positive;
this holds at the Wilson start and is preserved by the exact block
kernel. Smallness of signed activities alone is not a substitute
for this condition.
Use the same forest, path realization, clipped frozen source and
cutoff choices throughout the local integral construction of
Chapter~\ref{ch06:large-fields}. Construct the outgoing path by the
classical and quantum matching in Chapter~\ref{ch07:matching}.
Choose the stage constants, ratio floors, minimum
auxiliary periods and inverse coupling thresholds in the order
specified above.
Require also that the actual coarse period exceeds the largest
local read diameter and twice the largest auxiliary reflection
half-period used in the construction. This includes the minimum
period $100(B_m+1)$ in Chapter~\ref{ch06:large-fields}; enlarging
the stage-dependent period floor ensures that every localization
cube and its dependence neighborhood embeds in the actual torus.

Then every admissible continuous input curve has a
continuous output representation with fixed labels,
satisfying the regular and large-field bounds
\eqref{ch04:activity-bounds} and the pointwise identity
\[
 \operatorname{dens}(\rho_u)
       =\mathcal B_m^{\beta'_u}\operatorname{dens}(\rho'_u).
\]
Its inverse coupling is \eqref{ch09:reset} and obeys \eqref{ch09:drift}, uniformly in
volume, cutoff and curve parameter. The construction
retains restriction compatibility and all exact-cover
cutoff conditions. Output inverse coupling-band membership additionally
requires the numerical band margin in the choice of
scales.
\end{theorem}

\begin{proof}
The proof consists of the convergent local expansion,
the three remainder estimates, the exact identities
\eqref{ch09:subtract-identity} and
\eqref{ch09:path-identity}, and
Proposition~\ref{ch09:fusion-prop}. Positivity and
normalization refer to the actual block integral.
\end{proof}

This is a theorem for one compatible adjacent pair. Its premises
include the actual quadratic coefficient and fixed path coefficient bounds;
they cannot be replaced by unrelated bounds of the same numerical
size. The common Gaussian estimate is supplied by the complete
factor construction in Chapter~\ref{ch06:large-fields}, and
Sections~\ref{ch09:stationary-remainder}--\ref{ch09:centered-log}
give the auxiliary estimate for the same pure integral. The
classical bound contraction and inherited-loop contraction of
Chapter~\ref{ch07:matching} allow the path bounds to be retained
with the stated order of choices.

What remains for a whole trajectory is to make those choices
simultaneously, tune the Wilson inverse coupling to the prescribed coarse
endpoint, and prove stability as the ultraviolet cutoff changes.
These are additional assertions, not consequences of applying an
arbitrary finite step repeatedly. In particular, no higher parameter
analyticity is inferred from the continuity proved here.
